% p3-13-phd-integration

\section{P3 §13. PhD Program Integration}

13. PhD Program Integration

   Figure (fig-p3-ch13-program-integration). Programme-integration triptych --- Mapping table / 42
                                        catalogue / Coq map.

  This section maps the elements of the Trinity $S^{3}$AI / Flos Aureus PhD programme to the
  Paper 3 formalism. The purpose is to make explicit the cross-references between the
  two bodies of work, facilitating readers who arrive from the PhD monograph context.

  13.1 Mapping Summary

  The PhD-to-Paper-3 mapping is intentionally narrow:

  - Trinity anchor identity: from PhD Chapter 3 / flos37.tex to Equation (star), Sections
  2.1, 2.4, and 5.5; role: root algebraic invariant forcing the factor 3 in Trinity monomials.
  - Lucas ring: from the PhD Trinity Identity chapter to Definition 2.3, Corollary 4.5, and
  Proposition 5.9; role: discrete state space and exact termination locus.
  - 42 precision catalogue: from PhD Chapter 34 to Definition 7.3 and Algorithm 9.5; role:
  empirical target set for the catalogue MDL vector M, not proof.
  - Pellis expansion: from companion Papers 1 and 2 to Section 4 and Theorem 4.4; role:
  fine-scale correction and interference layer.
  - Trinity monomial: from companion Paper 2 to Section 3 and Definition 2.2; role:
  coarse-grained symbolic projection.
  - Algebraic AlphaPhi: from PhD Chapter 38 to Section 10.5 and equation (2); role: $\varphi$-3/2
  approx 0.118, algebraic, terminating.
  - Logarithmic AlphaPhi: from PhD Chapter 38 to Section 10.5 and equation (3); role:
  2ln$\varphi$/pi approx 0.30635, transcendental, non-terminating.
  - Coq citation map: from PhD Appendix F to Section 11.9 and Appendix E; role: formal
  provenance only, not empirical validation.
  - Falsification ledger: from PhD Appendix B to Section 9.4 and Appendix D; role:
  acceptance and rejection protocol.
  - Sanctioned and forbidden seeds: from PhD Chapter 34 to Section 2.2 and Algorithms
  9.1--9.5; role: prevent hidden RNG cherry-picking.

  13.2 The Trinity Identity as Algebraic Root
  The identity $\varphi$2 + $\varphi$-2 = 3 is not merely symbolic decoration in the Paper 3 framework;
  it is the algebraic reason why 3 appears as a monomial generator in T. Specifically:

  - The Lucas ring Z[$\varphi$] has norm form N(a + b$\varphi$) = a2 + ab - b2. The element 3 has norm
  N(3) = 9; the identity $\varphi$2 + $\varphi$-2 = 3 expresses 3 as a self-conjugate sum within the ring.
  - In the transfer operator, the canonical weight w = $\varphi$-|c| produces a geometric decay
  governed by $\varphi$; the pressure function achieves a critical inverse temperature at beta =
  log$\varphi$, which by (star) corresponds to "energy level 3" in the thermodynamic analogy.
  - The complexity shell growth |SL| $\sim$ A $\cdot$ 2cL with c approx log2 3 reflects the 3-fold
  branching of the Lucas recurrence, itself a consequence of the Trinity anchor.

  13.3 The 42 Catalogue Under the MDL Lens
  The 42 precision catalogue was assembled as an empirical survey of notable numerical
  fits between constants and symbolic expressions. In the Paper 3 MDL framework, each
  entry is reinterpreted as follows:

  - Low-MDL entries (MDL ll L decimal): the symbolic form genuinely compresses the
  constant; necessary (but not sufficient) condition for the form to be meaningful.
  - High-MDL entries (MDL > L decimal): the proposed symbolic form adds no
  compression; the entry should be re-evaluated or removed.

  - Marginal entries (MDL approx L decimal): additional theory needed to determine
  whether the compression is structural.

  The MDL analysis is purely numerical; it does not distinguish a coincidence from a
  theoretical relationship.

  13.4 Coq Provenance Map: Role in Paper 3
  The paper uses a local, reviewer-safe status taxonomy rather than reproducing a global
  proof ledger:

  - Trinity identity $\varphi$2+$\varphi$-2=3: status Qed; used in Equation (star), Sections 2.1, 5.5, and
  13.2.
  - Lucas closure of Z[$\varphi$]: status Qed; used in Proposition 2.4, Corollary 4.5, and
  Proposition 5.9.
  - Fibonacci recurrence Fn+2=Fn+1+Fn: status Qed; used implicitly in Section 2.2 and
  in the Binet-formula discussion.
  - Lucas formula Ln=$\varphi$n+psin: status Qed; used in Definition 2.3.
  - Projection map injectivity:     not   yet   formalised;   conditional   on   the   working
  independence hypothesis.
  - MDL optimality of PiL: not yet formalised in Coq; stated as a paper theorem and
  marked as future formalisation work.
  - Pellis convergence bound: not yet formalised in Coq; used as an analytic result of this
  paper.
  - Spectral gap and phase-transition statements: open conjectures, not theorem claims.

  Exact global Coq counts are not reproduced in the article. They must be regenerated
  from the repository at submission time, and any theorem whose dependency chain
  touches Admitted material must be labelled conditional.

  13.5 Falsification Ledger Integration
  The falsification protocol specifies acceptance and rejection criteria for the overall PhD
  programme. Paper 3 contributes the following pre-registered tests (see also Section
  9.4):

  Test P3-T1 (MDL improvement). For at least 30 of the 42 catalogue entries, MDL(T\textasciicircum{}*, c
  res) < L decimal. Failure mode: fewer than 30 entries satisfy this --- Trinity monomial
  representation fails to provide systematic compression.

  Test P3-T2 (Phase transition exponent range). For at least 80\% of the 42 catalogue
  entries, muT $\in$ [0, 2]. Failure mode: > 20\% of entries have muT > 2 or diverge.

  Test P3-T3 (Lucas ring termination; regression). For each of the sanctioned
  Fibonacci/Lucas seeds, Algorithm 4.1 terminates in at most 2max(|a|,|b|) steps
  (Theorem). This is a software regression test, not a conjecture.

  Test P3-T4 (AlphaPhi disambiguation; confirmed). |alpha$\varphi$( alg) - alpha$\varphi$(log)| > 0.1.
  Confirmed numerically by the reference implementation.

  Test P3-T5 (Trinity anchor identity; confirmed). |$\varphi$2 + $\varphi$-2 - 3| < 10-48 at 50-digit
  precision. Confirmed numerically (and formally by Coq Qed).

  Test P3-T6 (No physical derivation; editorial). Paper contains no sentence asserting a
  physical constant is derived from $\varphi$ or from the Trinity framework. Confirmed by author
  review.

  13.6 Reviewer-Hardening Gate Imported from the PhD Monograph
  The PhD monograph contributes a useful editorial discipline that is adopted here as a
  reviewer-facing gate rather than as additional scientific evidence. The gate has four
  parts:

  1. Notation lock. Every overloaded symbol must be split before review. This is why
  alpha$\varphi$( alg)=$\varphi$-3/2 and alpha$\varphi$(log)=2ln$\varphi$/pi are treated as different objects, with
  different algebraic status and different failure modes.
  2. Constants lock. Public constants are cited against their current authoritative tables.
  If a measured value changes, the symbolic approximation is recomputed rather than
  silently preserved.
  3. Falsification lock. Every empirical or statistical claim must have a failure criterion.
  Claims that have no falsifier are editorially demoted to motivation, analogy, or
  conjecture.
  4. Proof-status lock. Coq-backed identities are cited with their proof status. Qed
  statements may be used as formal provenance; Admitted statements are proof sketches
  only; no silent Admitted-to-Qed upgrade is allowed.

  This gate is deliberately conservative. It strengthens the trilogy by making the most
  likely reviewer criticisms explicit: look-elsewhere effects, stale constants, notation drift,
  overclaiming physical derivation, and formal-proof overreach.

  13.7 Scott Olsen Golden-Balance Motif
  The Scott Olsen material is imported only as historical-geometric context, not as
  evidence for any catalogue entry. Olsen's The Golden Section: Nature's Greatest Secret
  presents the golden section as a cross-disciplinary motif spanning mathematics, natural
  form, art, architecture, music, and philosophy (Internet Archive bibliographic record).
  In Paper 3, this lineage is narrowed to one formal visual motif:

  $\varphi$-2, $\varphi$-1, 1, $\varphi$, $\varphi$2.

  This "golden balance" is used as a reciprocal-scale diagram around the unit element. It
  does not enter the MDL score, the null-model analysis, the Coq proof ledger, or the
  physical-mechanism discussion. Its role is editorial and pedagogical: it gives readers a
  stable visual language for the same reciprocal symmetry that appears algebraically in
  $\varphi$2+$\varphi$-2=3.

  13.8 Erratum Discipline and Claim-Tier Taxonomy
  The article adopts the erratum discipline of the PhD programme: if a claim fails, it is not
  patched rhetorically. It is either withdrawn, substituted by a narrower claim, or
  deferred to future work with an explicit failure mode.

  Four claim tiers govern the final text:

  - Tier A --- algebraic identity: exact statements such as $\varphi$2+$\varphi$-2=3. These may be
  treated as mathematical results.
  - Tier B --- symbolic compressibility: empirical or statistical claims about
  low-description-length representations. These require null models, MDL accounting,
  and out-of-sample tests.
  - Tier C --- physical mechanism candidate: proposed links from symbolic structure to
  physical theory. These remain conjectural unless embedded in QFT, RG flow, lattice
  theory, or another accepted physical framework.
  - Tier D --- historical, visual, or interpretive context: material such as Olsen's
  golden-section lineage and PhD-style visual architecture. These passages orient the
  reader but do not prove Tier B or Tier C claims.

  This taxonomy is the main protection against reviewer over-read. A beautiful diagram
  or historical analogy cannot upgrade a symbolic approximation into a physical
  derivation.

\subsection{References}

- Vasilev, D. \textit{Flos Aureus Monograph: Trinity Constants}. DOI \href{https://zenodo.org/doi/10.5281/zenodo.19227877}{10.5281/zenodo.19227877}.
- CODATA 2018 recommended values: NIST \href{https://physics.nist.gov/cuu/Constants/}{https://physics.nist.gov/cuu/Constants/}.
- Heyrovska, R. Golden-angle FSC publications --- Heyrovský Institute of Physical Chemistry.
